\section{Introduction}
\label{sec:intro}

Robotic manipulation requires understanding object affordances---the action possibilities
an object offers~\cite{myers2015affordance}.
Self-supervised video models, especially the JEPA family~\cite{bardes2025vjepa2},
transfer strongly to robotic tasks without task-specific fine-tuning.
Video SSL models like V-JEPA learn implicit world models through temporal prediction~\cite{lecun2022jepa}, making them natural candidates for embodied perception.
V-JEPA~2.1~\cite{bardes2026vjepa21} extends V-JEPA~2 with a Dense Predictive Loss applied
hierarchically at intermediate encoder layers, yielding a 20-point grasping improvement---yet
its evaluation considers only final-layer features, leaving internal representational structure unexamined.
DINOv2~\cite{oquab2024dinov2} provides our static-image SSL baseline.
Linear probes~\cite{alain2016probes} and attention head pruning~\cite{michel2019sixteen,elhage2021circuits}
are established interpretability tools; to our knowledge, we apply them to affordance features in frozen
video encoders for the first time.

We make three contributions:
(1) a layer-wise probe study (layers 0--24) showing video SSL models crystallize affordance
signal four layers earlier (layer~19) than DINOv2 (layer~23), with V-JEPA~2.1's intermediate
supervision producing stronger early-layer representations despite identical peak depth;
(2) an attention head ablation at each model's peak layer revealing sparse ``affordance heads''; and
(3) a pruning analysis showing ${\approx}5.3{\times}$ attention compute reduction at ${\leq}0.1\%$ mAP loss.
